% from ega3-1-fr.tex
\oldpage[III]{91}
\begin{lemma}[1.4.9]
\label{III.1.4.9-fr}
Soient $X$ un schéma quasi-compact,
$\mathfrak U=(U_i)_{1\leq i\leq n}$ un recouvrement fini de $X$ par des
ouverts affines, $\mathcal F$ un $\mathcal O_X$-Module quasi-cohérent. Le
complexe de faisceaux $\mathcal C^\bullet(\mathfrak U,\mathcal F)$ défini
par le recouvrement $\mathfrak U$ (G, II, 5.2) est alors un
$\mathcal O_X$-Module quasi-cohérent.
\end{lemma}

\begin{proof}
Il résulte des définitions (G, II, 5.2) que
$\mathcal C^p(\mathfrak U,\mathcal F)$ est somme directe des faisceaux
images directes des $\mathcal F\vert U_{i_0\ldots i_p}$ par l'injection
canonique $U_{i_0\ldots i_p}\to X$. Or, l'hypothèse que $X$ est un schéma
entraîne que ces injections sont des morphismes affines (I, 5.5.6), donc les
$\mathcal C^p(\mathfrak U,\mathcal F)$ sont quasi-cohérents (II, 1.2.6).
\end{proof}

\begin{proposition}[1.4.10]
\label{III.1.4.10-fr}
Soit $u:X\to Y$ un morphisme séparé et quasi-compact. Pour tout
$\mathcal O_X$-Module quasi-cohérent $\mathcal F$, les
$R^qu_*(\mathcal F)$ sont des $\mathcal O_Y$-Modules quasi-cohérents.
\end{proposition}

\begin{proof}
La question étant locale sur $Y$, on peut supposer $Y$ affine. Alors $X$ est
réunion finie d'ouverts affines $U_i$ ($1\leq i\leq n$)~; soit
$\mathfrak U$ le recouvrement $(U_i)$. En outre, comme $Y$ est un schéma,
il résulte de (I, 5.5.10) que pour tout ouvert affine $V\subset Y$,
l'injection canonique $u^{-1}(V)\to X$ est un morphisme affine~; on en
conclut ((\hyperref[III.1.4.1-fr]{1.4.1}) et (G, II, 5.2)) que l'on a un
isomorphisme canonique
\begin{equation}
\label{III.1.4.10.1-fr}
\tag{1.4.10.1}
H^\bullet(u^{-1}(V),\mathcal F)
\xrightarrow{\sim}H^\bullet(\Gamma(V,\mathcal K^\bullet))
\end{equation}
où l'on a posé
$\mathcal K^\bullet=u_*(\mathcal C^\bullet(\mathfrak U,\mathcal F))$. En
vertu de (\hyperref[III.1.4.9-fr]{1.4.9}) et (I, 9.2.2),
$\mathcal K^\bullet$ est un $\mathcal O_Y$-Module quasi-cohérent~; par
ailleurs, il constitue un \emph{complexe de faisceaux} puisqu'il en est
ainsi de $\mathcal C^\bullet(\mathfrak U,\mathcal F)$. Il résulte alors de la
définition de la cohomologie
$\mathcal H^\bullet(\mathcal K^\bullet)$ (G, II, 4.1) que cette dernière est
formée de $\mathcal O_Y$-Modules quasi-cohérents (I, 4.1.1). Comme (pour $V$
affine dans $Y$) le foncteur $\Gamma(V,\mathcal G)$ est exact en
$\mathcal G$ dans la catégorie des $\mathcal O_Y$-Modules quasi-cohérents,
on a (G, II, 4.1)
\begin{equation}
\label{III.1.4.10.2-fr}
\tag{1.4.10.2}
H^\bullet(\Gamma(V,\mathcal K^\bullet))
=\Gamma(V,\mathcal H^\bullet(\mathcal K^\bullet)).
\end{equation}
Notons enfin qu'il résulte de la définition de l'homomorphisme canonique
\[
H^\bullet(\mathfrak U,\mathcal F)\longrightarrow H^\bullet(X,\mathcal F)
\]
donnée dans (G, II, 5.2), que si $V'\subset V$ est un second ouvert affine
de $Y$, le diagramme
\[
\xymatrix{
H^\bullet(u^{-1}(V),\mathcal F)\ar[r]^{\sim}\ar[d]
&H^\bullet(\Gamma(V,\mathcal K^\bullet))\ar[d]\\
H^\bullet(u^{-1}(V'),\mathcal F)\ar[r]^{\sim}
&H^\bullet(\Gamma(V',\mathcal K^\bullet))
}
\]
est commutatif. On conclut donc de ce qui précède que les isomorphismes
(\hyperref[III.1.4.10.1-fr]{1.4.10.1}) définissent un isomorphisme de
$\mathcal O_Y$-Modules
\begin{equation}
\label{III.1.4.10.3-fr}
\tag{1.4.10.3}
R^\bullet u_*(\mathcal F)
\xrightarrow{\sim}\mathcal H^\bullet(\mathcal K^\bullet)
\end{equation}
